\chapter{存储器管理}

\section{存储器管理基本概念}
\concept{管理对象与目标}
\begin{points}
  \item 存储管理的对象是主存储器，即内存。
  \item 主存通常分为系统区和用户区：系统区存放内核程序与数据结构，用户区存放用户进程代码和数据。
  \item 目标：为用户提供方便、安全、充分大的存储空间，支持大型应用和系统程序运行。
\end{points}
\plain{把这个知识点放回本章主线理解：它回答的是 OS 为了管理资源、隐藏硬件、保证并发正确性而采用的某个对象、规则或步骤。}
\exam{常考选择/填空/简答，让你判断概念归属、比较相近概念，或按“定义-作用-机制-易错点”展开。}


\concept{存储器管理四个功能}
\begin{points}
  \item 分配与回收：为进程分配内存，进程结束或换出时回收。
  \item 抽象与映射：让进程看到连续的逻辑地址空间，并把逻辑地址转换成物理地址。
  \item 保护与共享：防止进程非法访问他人空间，同时允许受控共享。
  \item 存储扩充：在逻辑上提供比实际物理内存更大的存储空间。
\end{points}
\plain{把这个知识点放回本章主线理解：它回答的是 OS 为了管理资源、隐藏硬件、保证并发正确性而采用的某个对象、规则或步骤。}
\exam{常考选择/填空/简答，让你判断概念归属、比较相近概念，或按“定义-作用-机制-易错点”展开。}


\concept{逻辑地址与物理地址}
\begin{points}
  \item 逻辑地址是 CPU 运行的进程所看到的地址，也叫虚地址或相对地址。
  \item 物理地址是内存中的实际地址，也叫实地址或绝对地址。
  \item 地址变换/地址映射/重定位是把逻辑地址转换为物理地址的过程。
  \item MMU 是常见的硬件支持，负责地址转换和合法性检查。
\end{points}
\plain{程序看到的是逻辑地址，内存条上真正用的是物理地址，中间靠装载器/MMU 转换。}
\exam{常考逻辑地址到物理地址转换、静态/动态重定位、越界检查。}


\concept{地址绑定时机}
\begin{points}
  \item 编译时绑定：若装入地址已知，编译器可直接生成绝对地址；位置改变则需重新编译。
  \item 装入时绑定：编译生成可重定位代码，装入时确定物理地址。
  \item 运行时绑定：进程运行中仍可移动，地址在执行时由硬件转换，现代分页/分段常用。
\end{points}
\plain{程序看到的是逻辑地址，内存条上真正用的是物理地址，中间靠装载器/MMU 转换。}
\exam{常考逻辑地址到物理地址转换、静态/动态重定位、越界检查。}


\concept{静态重定位与动态重定位}
\begin{points}
  \item 静态重定位在程序装入内存时一次性完成地址修改，运行中不能随意移动。
  \item 动态重定位在程序运行时进行，每次访问内存时由重定位寄存器或 MMU 转换。
  \item 动态重定位支持进程换入换出、移动和虚拟存储，但需要硬件支持。
\end{points}
\plain{程序看到的是逻辑地址，内存条上真正用的是物理地址，中间靠装载器/MMU 转换。}
\exam{常考逻辑地址到物理地址转换、静态/动态重定位、越界检查。}


\section{连续分配管理}
\concept{单一连续分配}
\begin{points}
  \item 内存用户区一次只装入一道用户程序，用户程序独占用户区。
  \item 实现简单，适合早期单道系统。
  \item 缺点是内存利用率低，不能支持多道程序并发。
\end{points}
\plain{连续分配要求一整块内存，容易产生碎片；分页把连续要求打碎。}
\exam{常考内部/外部碎片区别、首次/最佳/最坏适应算法、伙伴系统合并。}


\concept{固定分区分配}
\begin{points}
  \item 系统启动时把内存划分为若干固定大小分区，每个分区装入一个进程。
  \item 分区大小可相等，也可不等。
  \item 优点：实现简单，支持多道程序。
  \item 缺点：分区内部未用空间形成内部碎片；分区数限制多道程序道数。
\end{points}
\plain{连续分配要求一整块内存，容易产生碎片；分页把连续要求打碎。}
\exam{常考内部/外部碎片区别、首次/最佳/最坏适应算法、伙伴系统合并。}


\concept{动态分区分配}
\begin{points}
  \item 不预先固定分区，进程装入时根据需要动态划分连续内存区。
  \item 优点：比固定分区灵活，内部碎片少。
  \item 缺点：随着分配和回收，会产生外部碎片，需要紧凑/拼接。
\end{points}
\plain{连续分配要求一整块内存，容易产生碎片；分页把连续要求打碎。}
\exam{常考内部/外部碎片区别、首次/最佳/最坏适应算法、伙伴系统合并。}


\concept{内部碎片与外部碎片}
\begin{points}
  \item 内部碎片：已经分配给进程但进程用不完的空间，位于已分配区域内部。
  \item 外部碎片：空闲空间总量足够，但分散在各处，无法形成一个足够大的连续区域。
  \item 固定分区和分页容易有内部碎片；动态分区和分段容易有外部碎片。
  \item 直觉：内部碎片是“房间给你了但你没住满”；外部碎片是“空房很多但都太零散，放不下大团队”。
\end{points}
\plain{连续分配要求一整块内存，容易产生碎片；分页把连续要求打碎。}
\exam{常考内部/外部碎片区别、首次/最佳/最坏适应算法、伙伴系统合并。}


\concept{动态分区分配算法}
\begin{points}
  \item 首次适应 FF：从低地址开始找到第一个足够大的空闲分区。速度较快，但低地址碎片较多。
  \item 循环首次适应 NF：从上次查找结束位置继续查找，分布较均匀。
  \item 最佳适应 BF：选择能满足要求的最小空闲分区，尽量减少剩余空间，但容易产生很多小碎片。
  \item 最坏适应 WF：选择最大的空闲分区，剩余空间仍较大，但可能破坏大空闲区。
\end{points}
\plain{连续分配要求一整块内存，容易产生碎片；分页把连续要求打碎。}
\exam{常考内部/外部碎片区别、首次/最佳/最坏适应算法、伙伴系统合并。}


\concept{伙伴系统}
\begin{points}
  \item 伙伴系统按 $2^k$ 大小管理内存块，申请时分裂大块，释放时若相邻伙伴空闲则合并。
  \item 优点：分裂和合并规则简单，管理效率较高。
  \item 缺点：申请大小不一定正好是 $2^k$，可能产生内部碎片。
\end{points}
\plain{连续分配要求一整块内存，容易产生碎片；分页把连续要求打碎。}
\exam{常考内部/外部碎片区别、首次/最佳/最坏适应算法、伙伴系统合并。}


\section{分页存储管理}
\concept{分页基本思想}
\begin{points}
  \item 将逻辑地址空间划分为大小相等的页 page。
  \item 将物理内存划分为同样大小的页框/页帧 frame。
  \item 以页为单位分配内存，逻辑上连续，物理上可以分散。
  \item 分页取消了作业对物理连续内存的要求，避免外部碎片，但最后一页可能有内部碎片。
\end{points}
\plain{分页就是逻辑上按页切、物理上按页框放，页表负责记录页到页框的映射。}
\exam{常考页号/页内偏移拆分、有效访问时间、多级页表节省页表空间。}


\concept{页表}
\begin{points}
  \item 系统为每个进程建立一张页表，记录逻辑页号到物理块号的映射。
  \item 每个页面对应一个页表项，页表项通常包含物理块号和标志位。
  \item 标志位可能包括存在位、修改位、引用位、保护位等。
  \item 进程切换时，页表基址和长度等信息会被装入页表寄存器。
\end{points}
\plain{分页就是逻辑上按页切、物理上按页框放，页表负责记录页到页框的映射。}
\exam{常考页号/页内偏移拆分、有效访问时间、多级页表节省页表空间。}


\concept{分页地址结构与转换}
\begin{points}
  \item 逻辑地址由页号 $p$ 和页内偏移 $d$ 组成。
  \item 页大小为 $2^n$ 字节时，低 $n$ 位是页内偏移，高位是页号。
  \item 地址转换：用页号查页表得到物理块号 $f$，物理地址为 $f\times$ 页大小 $+d$。
  \item 若页表项无效或越界，则产生地址越界或缺页异常。
\end{points}
\plain{进程不是一直运行，它会在就绪、运行、阻塞等状态之间因调度或等待事件而移动。}
\exam{常考三/五/七状态转换图，尤其运行到阻塞、阻塞到就绪、就绪到运行的触发条件。}

\concept{快表 TLB 与有效访问时间}
\begin{points}
  \item 快表 TLB 是页表项的高速缓存，用于减少访问页表的内存次数。
  \item 无快表时，访问一次数据通常需先访问页表，再访问数据，共两次内存访问。
  \item 有快表时，若命中，直接得到物理块号，只需访问数据；若未命中，需要访问页表再访问数据。
  \item 有效访问时间公式常写为：$EAT=\alpha(t_{TLB}+t_M)+(1-\alpha)(t_{TLB}+2t_M)$，其中 $\alpha$ 为命中率。
\end{points}
\plain{分页就是逻辑上按页切、物理上按页框放，页表负责记录页到页框的映射。}
\exam{常考页号/页内偏移拆分、有效访问时间、多级页表节省页表空间。}


\concept{多级页表}
\begin{points}
  \item 多级页表把页表本身分页，只为实际用到的页表页分配内存。
  \item 解决页表过大、必须连续存放的问题。
  \item 缺点：地址转换需要多次查表，若无 TLB 支持会增加访问开销。
  \item 做题时根据每级索引位数判断每级表项数，根据页大小和表项大小判断一页能装多少表项。
\end{points}
\plain{分页就是逻辑上按页切、物理上按页框放，页表负责记录页到页框的映射。}
\exam{常考页号/页内偏移拆分、有效访问时间、多级页表节省页表空间。}


\section{分段与段页式存储管理}
\concept{分段基本思想}
\begin{points}
  \item 分段按程序的逻辑结构划分地址空间，如代码段、数据段、栈段、共享库段。
  \item 逻辑地址由段号和段内偏移组成。
  \item 段表项包含段基址、段长和保护信息。
  \item 地址转换时先检查段号和段内偏移是否合法，再由段基址加偏移得到物理地址。
\end{points}
\plain{分页按固定大小切，方便管理；分段按程序逻辑模块切，方便共享保护。段页式是先分段再分页。}
\exam{常考分页与分段比较、段号+段内地址合法性、段页式地址结构。}


\concept{分页与分段比较}
\begin{longtable}{p{0.20\textwidth}p{0.36\textwidth}p{0.36\textwidth}}
\toprule
比较项 & 分页 & 分段 \\
\midrule
划分依据 & 系统按固定大小机械划分 & 用户程序逻辑结构划分 \\
单位大小 & 页大小固定 & 段大小可变 \\
地址空间 & 一维地址空间 & 二维地址空间：段号+段内偏移 \\
碎片 & 无外部碎片，有内部碎片 & 无内部碎片或较少，有外部碎片 \\
共享保护 & 不如分段自然 & 以逻辑段为单位，便于共享与保护 \\
\bottomrule
\end{longtable}

\concept{段页式管理}
\begin{points}
  \item 段页式先把程序按逻辑分段，再把每个段分页。
  \item 地址结构通常为段号、段内页号、页内偏移。
  \item 每个段有段表项，段表项指向该段对应的页表。
  \item 优点：兼具分段的逻辑保护共享和分页的离散分配；缺点是地址转换更复杂。
\end{points}
\plain{分页按固定大小切，方便管理；分段按程序逻辑模块切，方便共享保护。段页式是先分段再分页。}
\exam{常考分页与分段比较、段号+段内地址合法性、段页式地址结构。}


